\chapter{Confronto finale e scelta per ruolo}
Il confronto non individua un controllore che soddisfi contemporaneamente ogni obiettivo. Le famiglie mettono in evidenza compromessi diversi, che diventano più chiari quando massa e ritardi variano insieme.

LQG stabilizza in modo robusto entrambe le rappresentazioni e limita bene il comando prodotto dal rumore. Il prefiltro offre precisione sul nominale, ma non corregge automaticamente l'offset dovuto a un carico diverso o a una coppia costante. È una soluzione utile per mostrare il ruolo dello stato stimato e il limite di un inseguimento senza integrazione.

LQGI aggiunge soltanto due stati e mantiene la stabilità robusta. L'integratore corregge l'errore misurato nei casi stabili e rende il tracking meno sensibile alla massa. Il prezzo osservato è il tempo di assestamento, che supera il target scelto. Per un compito di posizionamento lento e preciso è il candidato più comprensibile da approfondire, senza dichiararlo già conforme a tutte le specifiche.

Hinf1 mostra come la sensibilità mista controlli trasferimenti pesati senza garantire una buona forma del transitorio. Passare alla struttura 2-DoF permette a Hinf2 di seguire molto meglio il modello desiderato sul nominale. Tuttavia il suo margine robusto è piccolo e il rumore genera coppie molto elevate. La piccola norma nominale è quindi un risultato matematico corretto, ma insufficiente a motivarne l'uso.

Il PID strutturato contiene pochi stati e conserva una legge facilmente riconoscibile. Ottiene NP sul costo comune, ma non RS. La semplicità della struttura non elimina l'accoppiamento del braccio né il ritardo dei drive: i due canali diagonali hanno una libertà limitata per affrontarli.

La D--K incorpora direttamente la struttura incerta, ma il problema non è convesso e gli scaling aumentano l'ordine. Le due sintesi eseguite con uguali pesi e inizializzazione non raggiungono l'obiettivo robusto unitario. Mu\_A lascia un errore DC di circa il 99.8\%, mentre Mu\_B lascia circa il 21.9\% e amplifica il rumore fino a un picco osservato di circa 762 N m. Entrambi mantengono RS, ma queste prestazioni non ne motivano l’adozione. Il confronto rimane utile perché mostra quanto una descrizione più ricca o più libera possa cambiare la ricerca. Non scegliamo un candidato soltanto perché proviene da una sintesi robusta.

\section{Che cosa offre la rappresentazione parametrica}
A conserva il significato di ogni variazione. Una massa più grande modifica insieme inerzia e gravità; un ritardo più grande sposta insieme il polo e lo zero di Padé. Questa correlazione evita che l'analisi combini coefficienti fisicamente incompatibili. È la descrizione naturale quando gli intervalli dei parametri sono noti e il modello fisico è credibile.

Il costo è la complessità della struttura reale ripetuta: sette grandezze non significano sette sole porte. Prodotti, inversi e ripetizioni possono rendere più impegnative l'analisi e la sintesi. La fedeltà della famiglia non rende automaticamente semplice il calcolo.

\section{Che cosa offre la rappresentazione concentrata}
B riassume gli scostamenti dei drive in due perturbazioni dinamiche e un peso comune. Il procedimento SISO rende leggibile il passaggio dalla famiglia all'inviluppo e al fit. Il carico resta esplicito nella meccanica: il confronto riguarda davvero il modo di rappresentare gli attuatori.

Il prezzo è la libertà concessa a dinamiche che non corrispondono necessariamente a frequenze, smorzamenti e ritardi fisici della famiglia. Questa libertà può ridurre i margini dimostrabili. Inoltre il peso è verificato su un insieme finito di casi: la copertura globale resta un limite dichiarato.

Le due rappresentazioni sono quindi strumenti complementari. A aiuta a spiegare l'origine fisica delle variazioni; B offre una descrizione dinamica compatta dei drive coperti. Nei dieci impianti ricostruiti, a controllore fissato, producono la stessa risposta. Sono le garanzie sull'insieme ammesso a cambiare.
